\clearpage
\VerdierPhysicalPageAnchor{0179}%
\begingroup\small\setlength{\parskip}{0pt}\setlength{\emergencystretch}{1.5em}
\setlength{\abovedisplayskip}{4pt plus 2pt minus 2pt}%
\setlength{\belowdisplayskip}{4pt plus 2pt minus 2pt}%

\VerdierSemanticAnchor{ch3d-1-2-4}
\noindent\textbf{1.2.4.}\quad
$D(\mathcal A)$ の部分圏
$K^{,b}(\mathcal A)/\operatorname{Ac}(\mathcal A)$（それぞれ
$K^{,+}(\mathcal A)/\operatorname{Ac}(\mathcal A)$、それぞれ
$K^{,-}(\mathcal A)/\operatorname{Ac}(\mathcal A)$）を
$D^b(\mathcal A)$（それぞれ $D^+(\mathcal A)$、それぞれ
$D^-(\mathcal A)$）と書き、\emph{有界対象の部分圏}（それぞれ
\emph{下に有界な対象の部分圏}、それぞれ\emph{上に有界な対象の部分圏}）
と呼ぶ。これらの有界性条件は、考えている複体のコホモロジー対象にのみ
課されることに注意すべきである。

\VerdierSemanticAnchor{ch3d-1-2-5}
\medskip\noindent\textbf{1.2.5.}\quad
$X$ と $Y$ を $D(\mathcal A)$ の二つの対象とする。群
$\operatorname{Hom}_{D(\mathcal A)}(X,Y)$ を
$\operatorname{Ext}_{\mathcal A}^{0}(X,Y)$ と書く。同様に、群
$\operatorname{Ext}_{\mathcal A}^{0}(X,Y[i])\simeq
\operatorname{Ext}_{\mathcal A}^{0}(X[-i],Y)$ を
$\operatorname{Ext}_{\mathcal A}^{i}(X,Y)$ と書く（第 I 章，1.3.8）。
この記法は後に正当化される。ただし、$D(\mathcal A)$ の任意の完全三角
$(X,X',X'',u,v,w)$（それぞれ $(Y,Y',Y'',u,v,w)$）が次の両側無限完全列を
与えることを思い出しておく：
\[
\cdots\longrightarrow \operatorname{Ext}_{\mathcal A}^{0}(X'',Y)
\longrightarrow \operatorname{Ext}_{\mathcal A}^{0}(X',Y)
\longrightarrow \operatorname{Ext}_{\mathcal A}^{0}(X,Y)
\longrightarrow \operatorname{Ext}_{\mathcal A}^{1}(X'',Y)
\longrightarrow\cdots
\]
（それぞれ
\[
\cdots\longrightarrow \operatorname{Ext}_{\mathcal A}^{0}(X,Y)
\longrightarrow \operatorname{Ext}_{\mathcal A}^{0}(X,Y')
\longrightarrow \operatorname{Ext}_{\mathcal A}^{0}(X,Y'')
\longrightarrow \operatorname{Ext}_{\mathcal A}^{1}(X,Y)
\longrightarrow\cdots )
\]
（第 II 章，1.2.1）。

\VerdierSemanticAnchor{ch3d-1-2-6}
\medskip\noindent\textbf{1.2.6.}\quad
$M$ を $\mathcal A$ の対象の集合とする。すべての次数でコホモロジー対象が
$M$ に属する複体によって定まる $D(\mathcal A)$ の充満部分圏を
$D_M(\mathcal A)$ と書く。充満部分圏
$D_M(\mathcal A)\cap D^*(\mathcal A)$ を $D_M^*(\mathcal A)$ と書く。
$M$ の対象によって定まる $\mathcal A$ の充満部分圏が $\mathcal A$ の
厚い部分圏であるとき（\textit{すなわち}、それがアーベル圏であり、包含関手が
完全であり、かつ $M$ の対象の集合が拡大について閉じているとき）、部分圏
$D_M^*(\mathcal A)$ は $D(\mathcal A)$ の充満三角部分圏である。

\VerdierSemanticAnchor{ch3d-1-2-7}
\medskip\noindent\textbf{1.2.7.}\quad
$F\colon\mathcal A\longrightarrow\mathcal B$ を二つのアーベル圏の間の完全関手
とし、$K(F)\colon K(\mathcal A)\longrightarrow K(\mathcal B)$ を $F$ の
複体のホモトピー圏への拡張とする（第 I 章，2.5.10）。関手 $K(F)$ は
$\mathcal A$ の非輪状複体を $\mathcal B$ の非輪状複体に送り、完全であり
（第 II 章，1.3.7）、したがって（第 II 章，2.2.6）完全関手
$D(F)\colon D(\mathcal A)\longrightarrow D(\mathcal B)$ を定める。関手
$D(F)$ は $D(\mathcal A)$ の有界（それぞれ下に有界、それぞれ上に有界）な
対象を $D(\mathcal B)$ の同じ種類の対象に送る。多くの場合、関手 $D(F)$
を表すのに濫用的な記法 $F\colon D(\mathcal A)\longrightarrow D(\mathcal B)$
を用いる。

\VerdierSemanticAnchor{ch3d-1-2-8}
\medskip\noindent\textbf{1.2.8.}\quad
標準同型（第 II 章，1.3.8）
$i\colon K(\mathcal A)^o\longrightarrow K(\mathcal A^o)$ は
$\operatorname{Ac}(\mathcal A)^o$ の対象を
$\operatorname{Ac}(\mathcal A^o)$ の対象に送り、したがって、やはり
$i\colon D(\mathcal A)^o\longrightarrow D(\mathcal A^o)$ と書く三角圏の
同型を定める。
\endgroup
